\chapter{Related Work}
\label{ch:related_work}

“Projection differences” predict how training in some data would change the model’s behaviour

“Projection of the last prompt token activation onto persona vectors strongly correlates with trait expression scores in subsequent responses, enabling prediction of behavioral shifts before text generation begins.”

A projection difference is how much a training dataset's responses differ from the base model's responses along a feature (e.g., a personality trait) direction. It can be done by projecting into the weights or the activation’s space. 
[2507.21509] Persona Vectors: Monitoring and Controlling Character Traits in Language Models It identifies directions in the model's activation space (persona vectors) underlying several traits, such as evil, sycophancy, and propensity to hallucinate, and finds that both intended and unintended personality changes after fine-tuning are strongly correlated with shifts along the relevant persona vectors. The monitoring mechanism is projection onto those activation-space directions.
 [2406.11717] Refusal in Language Models Is Mediated by a Single Direction and [2502.17420] The Geometry of Refusal in Large Language Models: Concept Cones and Representational Independence suggest that refusal behavior is mediated by a single direction or a “cone, respectively
[2605.04572] From Parameter Dynamics to Risk Scoring : Quantifying Sample-Level Safety Degradation in LLM Fine-tuning They measure “safety” and “harmful” as linear directions in the gradient of the weights with respect the data. It proposes SQSD, which quantifies each sample's fine-tuning risk by the projection gap of its induced parameter update along danger versus safety directions. Concretely, it computes the sample-induced parameter update via one-step gradient, projects this update onto danger and safety directions for each module, and aggregates the projection gap across all modules.

Note: We can’t use a weight-space projection because it needs the sample- or dataset-induced parameter update: at minimum a backward pass, and the clean multi-epoch definition is the actual cumulative update, i.e. we’d have to run part of the training we're trying to forecast. 


Degradation requires studying the data×model interaction, not data content alone
[2310.03693] Fine-tuning Aligned Language Models Compromises Safety, Even When Users Do Not Intend To! “simply fine-tuning with benign and commonly used datasets can also inadvertently degrade the safety alignment of LLMs”
Out of Tune: Fine-Tuning Foundation Models Leads to Unpredictable Safety Drift - Center for Democracy and Technology "no single fine-tuning choice or combination of choices consistently explained why models became safer in some cases and less safe in others"; magnitude of model change also failed as a predictor. 


Drift accumulates roughly linearly in parameter space (projections onto danger directions sum across updates), but the behavioral readout is a threshold function
[2605.04572] From Parameter Dynamics to Risk Scoring : Quantifying Sample-Level Safety Degradation in LLM Fine-tuning
Observation: This means that a predictor model might be easier to build than expected. Although learning the threshold might be complicated.
Note: SQSD does not tell you (a) how far you currently are from the basin boundary, or (b) when the accumulated trajectory will cross it. Predicting these seems difficult. 
